\documentclass[11pt]{article}
\usepackage{amsmath,amssymb,amsthm}

\newtheorem{theorem}{Theorem}
\newtheorem{lemma}[theorem]{Lemma}
\newtheorem{corollary}[theorem]{Corollary}
\newtheorem{remark}[theorem]{Remark}

\title{Commuting unitary extensions of a commuting pair of isometries}
\author{}
\date{}

\begin{document}
\maketitle

\begin{theorem}
\label{thm:main}
Let \(H\) be a complex Hilbert space and let \(T,T'\in B(H)\) be isometries
satisfying \(TT'=T'T\). Then there exist a Hilbert space \(K\supset H\) and
unitaries \(U,U'\in B(K)\) such that
\[
U\big|_H=T,\qquad
U'\big|_H=T',\qquad
UU'=U'U.
\]
In particular, \(H\) is invariant under both \(U\) and \(U'\).
\end{theorem}

\noindent
The proof is the classical inductive construction via minimal unitary extensions.
We first record the single-operator fact, then extend the second isometry, then
unitarize.

\bigskip
%% ============================================================
\section*{Preliminary: minimal unitary extension of one isometry}

\begin{lemma}[Wold]
\label{lem:wold}
Let \(V\in B(H)\) be an isometry. Then there is an orthogonal decomposition
\(H=H_u\oplus H_s\) reducing \(V\), with \(V\big|_{H_u}\) unitary and
\(V\big|_{H_s}\) a unilateral shift. Writing \(L=H_s\ominus VH_s=\ker(V^*\big|_{H_s})\)
and identifying \(H_s\cong\bigoplus_{n\ge 0}L\), the bilateral shift \(\widetilde{S}\) on
\(K_s=\bigoplus_{n\in\mathbb{Z}}L\) is a unitary extension of \(V\big|_{H_s}\).
Consequently
\[
U_V:=V\big|_{H_u}\oplus\widetilde{S}
\quad\text{on}\quad
K_V:=H_u\oplus K_s
\]
is a unitary extension of \(V\). It is \emph{minimal} in the sense that
\begin{equation}
\label{eq:minimal}
K_V=\bigvee_{n\in\mathbb{Z}}U_V^n H.
\end{equation}
Moreover \(U_V^k\big|_H=V^k\) for every integer \(k\ge 0\).
\end{lemma}

\begin{remark}
The Halmos matrix
\(\displaystyle\begin{pmatrix}V&D_{V^*}\\ 0&-V^*\end{pmatrix}\)
on \(H\oplus H\) is always a unitary extension of an isometry \(V\), but it is
\emph{not} always minimal. The inductive construction below requires
minimality~\eqref{eq:minimal}, so we use Lemma~\ref{lem:wold}.
\end{remark}

\bigskip
%% ============================================================
\section*{Step 1 --- Minimal unitary extension of \(T\)}

Apply Lemma~\ref{lem:wold} to \(T\). Let \(U\) be a minimal unitary extension of
\(T\) on
\begin{equation}
\label{eq:Htilde}
\widetilde{H}
=
\bigvee_{n\in\mathbb{Z}}U^n H.
\end{equation}
Thus
\begin{equation}
\label{eq:U-powers}
U\big|_H=T,
\qquad
U^k\big|_H=T^k\quad(k\ge 0).
\end{equation}

\bigskip
%% ============================================================
\section*{Step 2 --- Isometric extension of \(T'\) commuting with \(U\)}

On the dense linear manifold of finite sums
\(\sum_{n\in\mathbb{Z}}U^n h_n\) with \(h_n\in H\), define
\begin{equation}
\label{eq:def-Ttilde}
\widetilde{T}'\Biggl(\sum_{n\in\mathbb{Z}}U^n h_n\Biggr)
=
\sum_{n\in\mathbb{Z}}U^n\,(T'h_n).
\end{equation}

\begin{lemma}
\label{lem:Ttilde}
The formula~\eqref{eq:def-Ttilde} is well-defined and extends by continuity to an
isometry \(\widetilde{T}'\) on all of \(\widetilde{H}\), satisfying
\begin{enumerate}
\item[\rm(i)] \(\widetilde{T}'\big|_H=T'\),
\item[\rm(ii)] \(\widetilde{T}'\,U=U\,\widetilde{T}'\),
\item[\rm(iii)] if \(T'\) is unitary, then \(\widetilde{T}'\) is unitary.
\end{enumerate}
\end{lemma}

\begin{proof}
It is enough to check that~\eqref{eq:def-Ttilde} preserves inner products on
finite sums. By sesquilinearity it suffices to compare
\(\langle U^n(T'h),\,U^m(T'g)\rangle\) with \(\langle U^n h,\,U^m g\rangle\) for
\(h,g\in H\) and \(n,m\in\mathbb{Z}\).

Assume first \(n\ge m\). Then, using that \(U\) is unitary,
\begin{align*}
\bigl\langle U^n(T'h),\,U^m(T'g)\bigr\rangle
&=
\bigl\langle U^{n-m}(T'h),\,T'g\bigr\rangle.
\end{align*}
Here \(n-m\ge 0\) and \(T'h\in H\), so~\eqref{eq:U-powers} gives
\(U^{n-m}(T'h)=T^{n-m}(T'h)\). Hence
\begin{align*}
\bigl\langle U^{n-m}(T'h),\,T'g\bigr\rangle
&=
\bigl\langle T^{n-m}T'h,\,T'g\bigr\rangle
=
\bigl\langle T'T^{n-m}h,\,T'g\bigr\rangle
\\
&=
\bigl\langle T^{n-m}h,\,g\bigr\rangle
=
\bigl\langle U^{n-m}h,\,g\bigr\rangle
=
\bigl\langle U^n h,\,U^m g\bigr\rangle,
\end{align*}
where the second equality used \(T^kT'=T'T^k\) (which follows by induction from
\(TT'=T'T\)), and the third used that \(T'\) is an isometry. The case \(m>n\) is
symmetric (or apply the same argument to the adjoint pairing).

Therefore~\eqref{eq:def-Ttilde} is well-defined and isometric on a dense
subspace, so it extends uniquely to an isometry \(\widetilde{T}'\) on
\(\widetilde{H}\).

(i) Take only the \(n=0\) term in~\eqref{eq:def-Ttilde}.

(ii) On finite sums,
\begin{align*}
\widetilde{T}'U\Biggl(\sum_n U^n h_n\Biggr)
&=
\widetilde{T}'\Biggl(\sum_n U^{n+1}h_n\Biggr)
=
\sum_n U^{n+1}(T'h_n)
=
U\widetilde{T}'\Biggl(\sum_n U^n h_n\Biggr).
\end{align*}
By density and continuity, \(\widetilde{T}'U=U\widetilde{T}'\) on
\(\widetilde{H}\).

(iii) If \(T'\) is unitary then \(\operatorname{Ran}T'=H\), so
\[
\operatorname{Ran}\widetilde{T}'
\supset
\bigvee_{n\in\mathbb{Z}}U^n(\operatorname{Ran}T')
=
\bigvee_{n\in\mathbb{Z}}U^n H
=
\widetilde{H}.
\]
An isometry with dense range is unitary.
\end{proof}

\bigskip
%% ============================================================
\section*{Step 3 --- Unitarize \(\widetilde{T}'\) and extend \(U\)}

If \(\widetilde{T}'\) is already unitary, set \(K=\widetilde{H}\) and \(U'=\widetilde{T}'\).
Then \((U,U')\) satisfies the claim of Theorem~\ref{thm:main}.

Now suppose \(\widetilde{T}'\) is not unitary. Apply Lemma~\ref{lem:wold} to the
isometry \(\widetilde{T}'\) on \(\widetilde{H}\): let \(U'\) be a minimal unitary
extension of \(\widetilde{T}'\) on
\begin{equation}
\label{eq:K}
K
=
\bigvee_{n\in\mathbb{Z}}(U')^n\widetilde{H}.
\end{equation}
Thus
\begin{equation}
\label{eq:Uprime-powers}
U'\big|_{\widetilde{H}}=\widetilde{T}',
\qquad
(U')^k\big|_{\widetilde{H}}=(\widetilde{T}')^k\quad(k\ge 0).
\end{equation}

Define an operator \(\widetilde{U}\) on the dense manifold of finite sums
\(\sum_{n\in\mathbb{Z}}(U')^n\xi_n\) with \(\xi_n\in\widetilde{H}\) by
\begin{equation}
\label{eq:def-Utilde}
\widetilde{U}\Biggl(\sum_{n\in\mathbb{Z}}(U')^n\xi_n\Biggr)
=
\sum_{n\in\mathbb{Z}}(U')^n\,(U\xi_n).
\end{equation}

\begin{lemma}
\label{lem:Utilde}
The formula~\eqref{eq:def-Utilde} is well-defined and extends to a unitary
operator \(\widetilde{U}\) on \(K\) such that
\begin{enumerate}
\item[\rm(i)] \(\widetilde{U}\big|_{\widetilde{H}}=U\) (hence \(\widetilde{U}\big|_H=T\)),
\item[\rm(ii)] \(\widetilde{U}\,U'=U'\,\widetilde{U}\).
\end{enumerate}
\end{lemma}

\begin{proof}
As in Lemma~\ref{lem:Ttilde}, compare inner products. For
\(\xi,\eta\in\widetilde{H}\) and \(n\ge m\),
\begin{align*}
\bigl\langle (U')^n(U\xi),\,(U')^m(U\eta)\bigr\rangle
&=
\bigl\langle (U')^{n-m}(U\xi),\,U\eta\bigr\rangle
=
\bigl\langle (\widetilde{T}')^{n-m}(U\xi),\,U\eta\bigr\rangle
\\
&=
\bigl\langle U(\widetilde{T}')^{n-m}\xi,\,U\eta\bigr\rangle
=
\bigl\langle (\widetilde{T}')^{n-m}\xi,\,\eta\bigr\rangle
=
\bigl\langle (U')^n\xi,\,(U')^m\eta\bigr\rangle,
\end{align*}
where we used~\eqref{eq:Uprime-powers}, the relation
\(U\widetilde{T}'=\widetilde{T}'U\) from Lemma~\ref{lem:Ttilde}(ii), and that \(U\) is
an isometry. The case \(m>n\) is symmetric. Thus~\eqref{eq:def-Utilde} defines an
isometry on a dense subspace of \(K\), and extends to an isometry
\(\widetilde{U}\) on \(K\).

(i) Take \(n=0\).

(ii) Immediate from~\eqref{eq:def-Utilde} on a dense set.

It remains to check that \(\widetilde{U}\) is surjective (hence unitary). Since \(U\) is
unitary on \(\widetilde{H}\), one has \(\operatorname{Ran}U=\widetilde{H}\). Therefore
\[
\operatorname{Ran}\widetilde{U}
\supset
\bigvee_{n\in\mathbb{Z}}(U')^n(\operatorname{Ran}U)
=
\bigvee_{n\in\mathbb{Z}}(U')^n\widetilde{H}
=
K.
\]
So \(\widetilde{U}\) is unitary.
\end{proof}

\bigskip
%% ============================================================
\section*{Conclusion}

\begin{proof}[Proof of Theorem~\ref{thm:main}]
Let \(U\) be a minimal unitary extension of \(T\) on \(\widetilde{H}\) as in
Step~1. Let \(\widetilde{T}'\) be the isometric extension of \(T'\) on
\(\widetilde{H}\) furnished by Lemma~\ref{lem:Ttilde}. If \(\widetilde{T}'\) is unitary,
take \(K=\widetilde{H}\) and \(U'=\widetilde{T}'\). Otherwise let \(U'\) and
\(\widetilde{U}\) be as in Lemma~\ref{lem:Utilde}. In either case one obtains a
Hilbert space \(K\supset H\) and commuting unitaries on \(K\) extending \(T\) and
\(T'\).
\end{proof}

\begin{remark}[minimality]
One may always replace \(K\) by the joint reducing subspace
\[
K_{\min}=\bigvee_{m,n\in\mathbb{Z}}U^m{U'}^n H
\]
to obtain a minimal commuting unitary extension.
\end{remark}

\begin{remark}[arbitrary families]
The same induction yields a commuting unitary extension for any family
(finite or infinite) of mutually commuting isometries: unitarize one at a time,
extending all previously unitarized operators by the formula~\eqref{eq:def-Ttilde}
at each stage.
\end{remark}

\bigskip
%% ============================================================
\section*{Explicit matrices when \(TT'\) is pure (BCL)}

Assume in addition that \(TT'\) is a pure isometry, i.e.\
\((TT')^{*n}\to 0\) strongly. Set
\[
\mathcal{W}=\ker(TT')^*=H\ominus TT'H,
\quad
W_1=\ker T^*,
\quad
W_2=\ker{T'}^*.
\]
Define a unitary \(U_{\mathcal{W}}\) on \(\mathcal{W}\) and a projection \(P\) by
\begin{equation}
\label{eq:UW}
U_{\mathcal{W}}
=
\begin{pmatrix}
T'\big|_{W_1} & 0 \\
0 & T^*\big|_{TW_2}
\end{pmatrix}
:
W_1\oplus TW_2
\;\longrightarrow\;
T'W_1\oplus W_2,
\qquad
P=P_{W_2}.
\end{equation}
On \(K=\ell^2(\mathbb{Z},\mathcal{W})\cong L^2(\mathbb{T};\mathcal{W})\) the operators
\begin{equation}
\label{eq:BCL}
\boxed{
\begin{aligned}
(U\xi)_n
&=
U_{\mathcal{W}}^*\,P\,\xi_n
+
U_{\mathcal{W}}^*\,P^\perp\,\xi_{n-1},
\\[1mm]
(U'\xi)_n
&=
P^\perp\,U_{\mathcal{W}}\,\xi_n
+
P\,U_{\mathcal{W}}\,\xi_{n-1}
\end{aligned}
}
\end{equation}
are commuting unitaries. Equivalently, \(U=M_{\Phi}\) and \(U'=M_{\Psi}\) where
\[
\Phi(z)=U_{\mathcal{W}}^*(P+zP^\perp),
\qquad
\Psi(z)=(P^\perp+zP)U_{\mathcal{W}}
\]
are unitary-valued for a.e.\ \(z\in\mathbb{T}\). The subspace
\(H\simeq H^2(\mathcal{W})=\bigoplus_{n\ge 0}\mathcal{W}\) is invariant for both, and
the restrictions are unitarily equivalent to \(T,T'\) (Berger--Coburn--Lebow).

This uses only commutativity of \(T\) and \(T'\), not double commutativity.
If \(TT'\) is not pure, split off the maximal unitary part
\(H_u=\bigcap_{n\ge 0}(TT')^n H\) (which reduces both \(T\) and \(T'\)) and apply
the above on the pure summand.

\end{document}
